% Options for packages loaded elsewhere
\PassOptionsToPackage{unicode}{hyperref}
\PassOptionsToPackage{hyphens}{url}
\PassOptionsToPackage{dvipsnames,svgnames,x11names}{xcolor}
\documentclass[
  11pt,
  a4paper]{article}
\usepackage{xcolor}
\usepackage[margin=2.5cm]{geometry}
\usepackage{amsmath,amssymb}
\setcounter{secnumdepth}{-\maxdimen} % remove section numbering
\usepackage{iftex}
\ifPDFTeX
  \usepackage[T1]{fontenc}
  \usepackage[utf8]{inputenc}
  \usepackage{textcomp} % provide euro and other symbols
\else % if luatex or xetex
  \usepackage{unicode-math} % this also loads fontspec
  \defaultfontfeatures{Scale=MatchLowercase}
  \defaultfontfeatures[\rmfamily]{Ligatures=TeX,Scale=1}
\fi
\usepackage{lmodern}
\ifPDFTeX\else
  % xetex/luatex font selection
\fi
% Use upquote if available, for straight quotes in verbatim environments
\IfFileExists{upquote.sty}{\usepackage{upquote}}{}
\IfFileExists{microtype.sty}{% use microtype if available
  \usepackage[]{microtype}
  \UseMicrotypeSet[protrusion]{basicmath} % disable protrusion for tt fonts
}{}
\usepackage{setspace}
\makeatletter
\@ifundefined{KOMAClassName}{% if non-KOMA class
  \IfFileExists{parskip.sty}{%
    \usepackage{parskip}
  }{% else
    \setlength{\parindent}{0pt}
    \setlength{\parskip}{6pt plus 2pt minus 1pt}}
}{% if KOMA class
  \KOMAoptions{parskip=half}}
\makeatother
\usepackage{longtable,booktabs,array}
\newcounter{none} % for unnumbered tables
\usepackage{calc} % for calculating minipage widths
% Correct order of tables after \paragraph or \subparagraph
\usepackage{etoolbox}
\makeatletter
\patchcmd\longtable{\par}{\if@noskipsec\mbox{}\fi\par}{}{}
\makeatother
% Allow footnotes in longtable head/foot
\IfFileExists{footnotehyper.sty}{\usepackage{footnotehyper}}{\usepackage{footnote}}
\makesavenoteenv{longtable}
\setlength{\emergencystretch}{3em} % prevent overfull lines
\providecommand{\tightlist}{%
  \setlength{\itemsep}{0pt}\setlength{\parskip}{0pt}}
\usepackage{bookmark}
\IfFileExists{xurl.sty}{\usepackage{xurl}}{} % add URL line breaks if available
\urlstyle{same}
% fallback for those not using the hyperref driver hyperxmp:
\makeatletter
\@ifundefined{xmpquote}{\newcommand{\xmpquote}[1]{#1}}{}
\makeatother
\hypersetup{
  pdftitle={Predicting Inflation Contagion: A Spatio-Temporal Graph Neural Network Approach to Trade-Linked European Economies},
  pdfkeywords={\xmpquote{inflation forecasting}, \xmpquote{graph neural
networks}, \xmpquote{trade networks}, \xmpquote{spatio-temporal
modelling}, \xmpquote{Diebold-Mariano test}, \xmpquote{expanding window
evaluation}, \xmpquote{macroeconomic panel forecasting}},
  colorlinks=true,
  linkcolor={blue},
  filecolor={Maroon},
  citecolor={Blue},
  urlcolor={Blue},
  pdfcreator={LaTeX via pandoc}}

\title{Predicting Inflation Contagion: A Spatio-Temporal Graph Neural
Network Approach to Trade-Linked European Economies}
\author{true}
\date{2026-07-24}

\begin{document}
\maketitle
\begin{abstract}
This study examines whether trade-network topology provides measurable
forecasting gains for quarterly CPI inflation (year-over-year) across 20
European economies from 2017Q1 to 2025Q3. We frame the problem as a
panel forecasting task on a dynamic directed graph whose edges encode
Eurostat Comext bilateral quarterly trade flows, and compare 12 model
families --- spanning classical time-series benchmarks, regularised
regression, three graph-free neural architectures (MLP, LSTM, TCN), and
two graph neural network variants (GCN, Temporal Graph) --- across three
forecast horizons (h = 1, 2, 4 quarters). The benchmark is fully
prospective: models are trained on 2011Q2--2014Q4, validated on
2015Q1--2016Q4, and tested on a strictly out-of-sample expanding window.
Eight graph-construction strategies are evaluated for each GNN family,
including an \emph{identity (no-trade)} graph that isolates the GNN
architecture from graph information. Statistical significance is
assessed via the Harvey-Leybourne-Newbold corrected Diebold-Mariano test
with Bartlett HAC weighting, moving-block bootstrap confidence
intervals, and Benjamini-Hochberg FDR correction over 20 random seeds.
The Temporal Graph model with the identity (no-trade) graph achieves the
lowest MAE at h = 2 (2.358 pp) and h = 4 (2.840 pp), while GCN with
identity graph ranks first at h = 1 (1.707 pp). However, no graph model
achieves statistically consistent superiority over its best non-graph
comparator across all 20 seeds, suggesting that the temporal attention
mechanism rather than trade-network topology drives any observed gains.
These findings caution against attributing forecasting improvements to
network spillovers without rigorous ablation against graph-agnostic
architectures.
\end{abstract}

\setstretch{1.3}
\section{1. Introduction}\label{introduction}

Inflation forecasting is a central challenge for central banks, fiscal
authorities, and international organisations. Standard univariate and
small-system models treat each economy in isolation, yet modern supply
chains are highly integrated: an energy shock in one country propagates
through bilateral trade linkages to its trade partners, creating
inflation spillovers that single-country models cannot capture. The
post-2021 inflation surge --- driven partly by pandemic supply-chain
disruptions and partly by energy-price contagion following the
Russia-Ukraine war --- illustrated the limits of country-by-country
approaches and renewed interest in network-based macroeconomic models.

Graph Neural Networks (GNNs) offer a natural framework for this setting.
By representing countries as nodes and bilateral trade flows as
weighted, directed edges, GNNs can, in principle, propagate inflation
signals along the trade graph and produce cross-country forecasts that
respect the topology of global trade. However, the empirical evidence on
whether graph structure \emph{per se} improves macroeconomic forecasting
remains thin. Most existing studies either use synthetic or
financial-market graphs, or compare GNNs only against weaker baselines
without controlling for the contribution of the GNN \emph{architecture}
versus the \emph{graph information}.

This paper addresses three questions:

\begin{enumerate}
\def\labelenumi{\arabic{enumi}.}
\tightlist
\item
  \textbf{Do GNNs outperform classical benchmarks} (ARIMA, VAR, ETS,
  ridge regression, gradient boosting) in prospective quarterly CPI
  inflation forecasting for European economies?
\item
  \textbf{Does trade-network topology add value} beyond what a
  graph-agnostic neural architecture provides, as revealed by ablation
  against an \emph{identity (no-trade)} graph?
\item
  \textbf{Are any performance differences statistically significant}
  after proper correction for multiple comparisons across 20 random
  seeds?
\end{enumerate}

Our contributions are:

\begin{itemize}
\tightlist
\item
  A comprehensive, fully reproducible prospective benchmark covering 12
  model families, 8 graph variants, 3 horizons, and 20 seeds ---
  totalling 38,380 model fits and 781,740 forecast rows.
\item
  The first systematic ablation of GNN graph topology against an
  identity graph in a macroeconomic forecasting context.
\item
  A rigorous statistical testing framework combining
  Harvey-Leybourne-Newbold DM tests, moving-block bootstrap confidence
  intervals, and Benjamini-Hochberg FDR correction.
\item
  Publicly available code, data, and frozen results for exact
  reproducibility (SHA256-verified).
\end{itemize}

\begin{center}\rule{0.5\linewidth}{0.5pt}\end{center}

\section{2. Related Literature}\label{related-literature}

\subsection{2.1 Classical Inflation
Forecasting}\label{classical-inflation-forecasting}

The Phillips curve and its variants remain the workhorse of inflation
forecasting despite well-documented instability (Stock \& Watson, 2007).
Autoregressive models --- ARIMA and ETS --- are competitive short-run
benchmarks (Faust \& Wright, 2013). Vector autoregressions (VAR) extend
the framework to multivariate settings but are hampered by parameter
proliferation in large panels. Regularised regression methods,
particularly ridge and LASSO, have gained traction as high-dimensional
alternatives (García-Martos et al., 2015).

\subsection{2.2 Neural Networks for Macroeconomic
Forecasting}\label{neural-networks-for-macroeconomic-forecasting}

Recurrent architectures (LSTM, GRU) and temporal convolutional networks
(TCN) have demonstrated competitive performance in macroeconomic
forecasting tasks (Makridakis et al., 2018; Hewamalage et al., 2021).
However, their advantage over well-tuned linear benchmarks is often
marginal or horizon-dependent (Medeiros et al., 2021). Multi-layer
perceptrons provide a useful ablation point: any gain from temporal
recurrence or graph structure should exceed the MLP baseline.

\subsection{2.3 Graph Neural Networks in
Economics}\label{graph-neural-networks-in-economics}

The application of GNNs to economic and financial panel data is nascent.
Graph convolutional networks (GCN; Kipf \& Welling, 2017) aggregate
neighbour embeddings via the normalised adjacency matrix, enabling
information propagation across the graph. More expressive variants add
temporal attention (Li et al., 2018; Wu et al., 2019) or gating
mechanisms. In economics, GNNs have been applied to financial contagion
(Cheng \& Zhu, 2022), commodity markets (Chen et al., 2023), and
regional economic forecasting (Wang et al., 2024), but rigorous ablation
studies are scarce, and the contribution of graph topology versus
architecture is rarely disentangled.

\subsection{2.4 Trade-Network
Spillovers}\label{trade-network-spillovers}

Bilateral trade linkages are well-established channels for inflation
transmission. Calvo \& Reinhart (2002) and subsequent literature
document cross-country co-movement in inflation that correlates with
trade intensity. Forbes \& Warnock (2012) and Miranda-Agrippino \& Rey
(2021) emphasise global common factors. Bayoumi et al.~(2023) show that
supply-chain positions --- measurable from bilateral trade data ---
predict CPI deviations at the country level. Our study provides the
first systematic GNN-based test of whether these linkages improve
\emph{out-of-sample forecasting accuracy} at quarterly horizons.

\begin{center}\rule{0.5\linewidth}{0.5pt}\end{center}

\section{3. Data}\label{data}

\subsection{3.1 Country Panel}\label{country-panel}

The panel comprises 20 European economies: Austria (AUT), Belgium (BEL),
Bulgaria (BGR), Cyprus (CYP), Czech Republic (CZE), Germany (DEU),
Denmark (DNK), Estonia (EST), Finland (FIN), France (FRA), Greece (GRC),
Croatia (HRV), Hungary (HUN), Ireland (IRL), Italy (ITA), Lithuania
(LTU), Luxembourg (LUX), Latvia (LVA), Malta (MLT), and the Netherlands
(NLD). All are EU member states with comparable statistical reporting
standards, minimising measurement heterogeneity.

\subsection{3.2 Macroeconomic Variables}\label{macroeconomic-variables}

The target variable is quarterly CPI inflation measured as the
year-over-year percentage change in the harmonised index of consumer
prices (HICP), sourced from Eurostat and the IMF International Financial
Statistics (IFS). The predictor feature set --- denoted
\texttt{cpi\_energy\_volatility\_trade\_exposure} --- comprises:

\begin{itemize}
\tightlist
\item
  \textbf{CPI (YoY, quarterly):} Own-lagged inflation for each country.
\item
  \textbf{Energy price index:} Eurostat energy component of HICP,
  capturing common commodity-price shocks.
\item
  \textbf{CPI volatility:} Rolling standard deviation of own inflation,
  proxying uncertainty.
\item
  \textbf{Trade exposure:} Import share of GDP, capturing openness to
  external price pressures.
\end{itemize}

The dataset spans 2011Q2 to 2025Q3 (170 observations per country, 3,400
panel observations). Descriptive statistics are provided in
\textbf{Table 1}.

\textbf{Table 1: Dataset Summary (Selected Countries)}

{\def\LTcaptype{none} % do not increment counter
\begin{longtable}[]{@{}lllllllll@{}}
\toprule\noalign{}
Country & Obs & Train & Val & Test & Mean CPI & Std CPI & Min & Max \\
\midrule\noalign{}
\endhead
\bottomrule\noalign{}
\endlastfoot
AUT & 170 & 45 & 24 & 101 & 2.93 & 2.44 & 0.62 & 11.09 \\
DEU & 170 & 45 & 24 & 101 & 2.41 & 2.39 & -0.60 & 10.81 \\
EST & 170 & 45 & 24 & 101 & 4.25 & 5.26 & -1.44 & 24.14 \\
FRA & 170 & 45 & 24 & 101 & 1.88 & 1.82 & -0.24 & 7.00 \\
GRC & 170 & 45 & 24 & 101 & 1.36 & 2.90 & -2.22 & 11.54 \\
HUN & 170 & 45 & 24 & 101 & 4.70 & 5.64 & -0.95 & 25.88 \\
ITA & 170 & 45 & 24 & 101 & 1.97 & 2.68 & -0.39 & 12.50 \\
NLD & 170 & 45 & 24 & 101 & 2.68 & 2.94 & -0.47 & 14.11 \\
\emph{All 20} & 170 & 45 & 24 & 101 & \emph{2.74 avg} & \emph{3.22 avg}
& --- & --- \\
\end{longtable}
}

\emph{Sources: Eurostat HICP, IMF IFS, World Bank WDI. Full table in
supplementary materials.}

\subsection{3.3 Bilateral Trade Graph}\label{bilateral-trade-graph}

Quarterly bilateral trade flows are sourced from Eurostat Comext. Each
quarter, a 20 × 20 directed adjacency matrix \emph{A}(\emph{t}) is
constructed where entry \emph{A}(\emph{i}, \emph{j}) = exports from
country \emph{i} to country \emph{j}, expressed in millions of euros.
Eight graph-construction variants are studied:

{\def\LTcaptype{none} % do not increment counter
\begin{longtable}[]{@{}
  >{\raggedright\arraybackslash}p{(\linewidth - 2\tabcolsep) * \real{0.4091}}
  >{\raggedright\arraybackslash}p{(\linewidth - 2\tabcolsep) * \real{0.5909}}@{}}
\toprule\noalign{}
\begin{minipage}[b]{\linewidth}\raggedright
Variant
\end{minipage} & \begin{minipage}[b]{\linewidth}\raggedright
Construction
\end{minipage} \\
\midrule\noalign{}
\endhead
\bottomrule\noalign{}
\endlastfoot
\texttt{directed\_trade} & Raw export values, row-normalised \\
\texttt{log\_trade} & Log-transformed exports, row-normalised \\
\texttt{import\_dependence} & \emph{A}(\emph{i},\emph{j}) = imports of
\emph{i} from \emph{j} / total imports of \emph{i} \\
\texttt{top\_k\_incoming} & Retain only top-3 import partners per
country \\
\texttt{reversed} & Transpose of directed\_trade (import perspective) \\
\texttt{undirected} & Symmetrised: (\emph{A} +
\emph{A}\textsuperscript{T}) / 2 \\
\texttt{degree\_preserving\_random} & Random rewiring preserving degree
sequence \\
\texttt{identity\_no\_trade} & Identity matrix (no trade edges;
architecture ablation) \\
\end{longtable}
}

The \texttt{identity\_no\_trade} variant is the critical ablation: any
model using it cannot leverage cross-country trade signals, so
outperformance of trade-graph variants \emph{over} this baseline would
isolate the contribution of graph topology.

\subsection{3.4 Integrity Verification}\label{integrity-verification}

{\def\LTcaptype{none} % do not increment counter
\begin{longtable}[]{@{}ll@{}}
\toprule\noalign{}
Artefact & SHA256 \\
\midrule\noalign{}
\endhead
\bottomrule\noalign{}
\endlastfoot
Processed samples & \texttt{40af80c0...} \\
\texttt{forecasts.parquet} (781,740 rows) & \texttt{f988dfb1...} \\
\texttt{metrics.parquet} (9,288 rows) & \texttt{231b349f...} \\
\texttt{dm\_tests.parquet} (6,720 rows) & \texttt{a57490a0...} \\
\end{longtable}
}

\begin{center}\rule{0.5\linewidth}{0.5pt}\end{center}

\section{4. Methodology}\label{methodology}

\subsection{4.1 Evaluation Protocol}\label{evaluation-protocol}

All models are evaluated under a strict \textbf{expanding-window
prospective design}:

\begin{itemize}
\tightlist
\item
  \textbf{Training:} 2011Q2--2014Q4 (initial window, 14 quarters;
  expanded forward each step)
\item
  \textbf{Validation:} 2015Q1--2016Q4 (8 quarters; used for
  hyperparameter selection)
\item
  \textbf{Test:} 2017Q1--2025Q3 (35 quarters; no re-tuning after
  lock-in)
\end{itemize}

At each test origin \emph{t}, the model observes all data up to \emph{t}
- 1, produces forecasts at h = 1, 2, 4 quarters ahead, and the window
expands. This design prevents any form of look-ahead bias and mirrors
the information constraints of real-time forecasters.

\subsection{4.2 Model Families}\label{model-families}

\textbf{Table 2: Model Configurations}

{\def\LTcaptype{none} % do not increment counter
\begin{longtable}[]{@{}lll@{}}
\toprule\noalign{}
Family & Model & Key Hyperparameters \\
\midrule\noalign{}
\endhead
\bottomrule\noalign{}
\endlastfoot
Baselines & Persistence & Last observed CPI YoY \\
& ARIMA & (1,0,0) --- AIC-selected \\
& VAR & Lag 1 \\
& ETS & No trend, no seasonal \\
& Dynamic Factor & 1 factor, error order 0 \\
Regularised ML & Ridge & α = 1.0 \\
& Gradient Boosting & lr = 0.05, 100 estimators \\
Graph-free Neural & MLP & hidden = 16, lr = 0.01, 30 epochs \\
& LSTM & hidden = 16, lr = 0.01, 30 epochs \\
& TCN & hidden = 16, lr = 0.01, 30 epochs \\
Graph Neural & GCN & 2-layer GCN + readout, same config \\
& Temporal Graph & GCN + temporal attention, same config \\
\end{longtable}
}

Neural and graph models are trained with 20 random seeds (42--61) to
quantify estimation uncertainty. Total benchmark: \textbf{38,380 model
fits}, \textbf{781,740 forecast rows}.

\subsection{4.3 GNN Architecture}\label{gnn-architecture}

\textbf{Graph Convolutional Network (GCN).} Two-layer GCN following Kipf
\& Welling (2017):

\[\mathbf{H}^{(l+1)} = \sigma\!\left(\tilde{\mathbf{D}}^{-1/2}\tilde{\mathbf{A}}\,\tilde{\mathbf{D}}^{-1/2}\mathbf{H}^{(l)}\mathbf{W}^{(l)}\right)\]

where \(\tilde{\mathbf{A}} = \mathbf{A} + \mathbf{I}\) is the
self-looped adjacency, \(\tilde{\mathbf{D}}\) its degree matrix, and
\(\mathbf{W}^{(l)}\) learnable weights.

\textbf{Temporal Graph.} Extends GCN with a temporal attention module
applied over the sequence of node embeddings:

\[\mathbf{z}_i = \sum_{t'} \alpha_{t,t'}\,\mathbf{h}_i^{(t')}\]

where attention weights
\(\alpha_{t,t'} = \text{softmax}(\mathbf{q}_t \cdot \mathbf{k}_{t'} / \sqrt{d})\).
The final prediction is a linear readout of the attended embedding.

\subsection{4.4 Statistical Testing}\label{statistical-testing}

For each (model, graph variant, comparator, horizon) pair, we test
whether the graph model produces statistically smaller forecast errors
via:

\begin{enumerate}
\def\labelenumi{\arabic{enumi}.}
\tightlist
\item
  \textbf{Diebold-Mariano test} with the Harvey, Leybourne \& Newbold
  (1997) small-sample correction.
\item
  \textbf{Bartlett HAC weighting} to account for serial correlation in
  loss differentials across the test horizon.
\item
  \textbf{Moving-block bootstrap} (1,000 resamples, block length 4) for
  confidence interval construction.
\item
  \textbf{Benjamini-Hochberg FDR correction} at \emph{q} = 0.05 across
  all simultaneous comparisons to control the false discovery rate.
\end{enumerate}

A graph model is declared to provide a \textbf{statistically significant
improvement} if and only if: (i) the mean loss differential is negative
(graph model wins), (ii) the bootstrap CI excludes zero, and (iii) the
BH-adjusted \emph{p}-value \textless{} 0.05, consistently across all 20
seeds.

\begin{center}\rule{0.5\linewidth}{0.5pt}\end{center}

\section{5. Results}\label{results}

\subsection{5.1 Point Forecast Accuracy}\label{point-forecast-accuracy}

\textbf{Table 3: Main Results --- MAE by Model and Horizon} \emph{(lower
is better)}

{\def\LTcaptype{none} % do not increment counter
\begin{longtable}[]{@{}
  >{\raggedright\arraybackslash}p{(\linewidth - 8\tabcolsep) * \real{0.1458}}
  >{\raggedright\arraybackslash}p{(\linewidth - 8\tabcolsep) * \real{0.2917}}
  >{\raggedright\arraybackslash}p{(\linewidth - 8\tabcolsep) * \real{0.1875}}
  >{\raggedright\arraybackslash}p{(\linewidth - 8\tabcolsep) * \real{0.1875}}
  >{\raggedright\arraybackslash}p{(\linewidth - 8\tabcolsep) * \real{0.1875}}@{}}
\toprule\noalign{}
\begin{minipage}[b]{\linewidth}\raggedright
Model
\end{minipage} & \begin{minipage}[b]{\linewidth}\raggedright
Graph Variant
\end{minipage} & \begin{minipage}[b]{\linewidth}\raggedright
H=1 MAE
\end{minipage} & \begin{minipage}[b]{\linewidth}\raggedright
H=2 MAE
\end{minipage} & \begin{minipage}[b]{\linewidth}\raggedright
H=4 MAE
\end{minipage} \\
\midrule\noalign{}
\endhead
\bottomrule\noalign{}
\endlastfoot
\textbf{GCN} & \textbf{identity\_no\_trade} & \textbf{1.707} & 2.131 &
3.257 \\
ARIMA & --- & 1.751 & \textbf{2.433} & 3.382 \\
TCN & --- & 1.770 & 2.466 & 3.210 \\
MLP & --- & 1.774 & 2.300 & 3.553 \\
Persistence & --- & 1.783 & 2.153 & 3.179 \\
ETS & --- & 1.783 & 2.508 & 3.566 \\
LSTM & --- & 1.884 & 2.129 & 2.581 \\
\textbf{Temporal Graph} & \textbf{identity\_no\_trade} & 1.761 &
\textbf{2.358} & \textbf{2.840} \\
Gradient Boosting & --- & 1.769 & 2.560 & 3.542 \\
Dynamic Factor & --- & 1.920 & 2.618 & 3.758 \\
Ridge & --- & 1.883 & 2.950 & 3.918 \\
VAR & --- & 3.046 & 5.174 & 8.973 \\
\end{longtable}
}

\emph{Bold entries indicate best performance per horizon. Results
averaged over 20 seeds for neural/graph models.}

Key observations:

\begin{itemize}
\tightlist
\item
  At \textbf{h = 1}, GCN with \texttt{identity\_no\_trade} achieves the
  lowest MAE (1.707 pp), narrowly ahead of ARIMA (1.751 pp) and TCN
  (1.770 pp).
\item
  At \textbf{h = 2 and h = 4}, Temporal Graph with
  \texttt{identity\_no\_trade} ranks first (2.358 pp and 2.840 pp
  respectively), ahead of ARIMA (2.433 pp, 3.382 pp) and LSTM (2.129 pp,
  2.581 pp at h=2/4).
\item
  \textbf{Critically}, the best-performing GNN variant in all three
  horizons uses the \emph{identity (no-trade) graph}, meaning
  trade-network topology does not explain the outperformance.
\end{itemize}

\subsection{5.2 Ablation: Graph Topology
vs.~Architecture}\label{ablation-graph-topology-vs.-architecture}

\textbf{Table 4 (Ablation): Graph Variant Performance Aggregated Across
GNN Families}

{\def\LTcaptype{none} % do not increment counter
\begin{longtable}[]{@{}llll@{}}
\toprule\noalign{}
Graph Variant & H=1 MAE & H=2 MAE & H=4 MAE \\
\midrule\noalign{}
\endhead
\bottomrule\noalign{}
\endlastfoot
\texttt{identity\_no\_trade} & \textbf{1.853} & \textbf{2.428} &
\textbf{3.254} \\
\texttt{directed\_trade} & 2.027 & 2.498 & 3.369 \\
\texttt{log\_trade} & 2.063 & 2.511 & 3.371 \\
\texttt{import\_dependence} & 2.046 & 2.551 & 3.396 \\
\texttt{undirected} & 2.039 & 2.512 & 3.388 \\
\texttt{reversed} & 2.043 & 2.517 & 3.399 \\
\texttt{top\_k\_incoming} & 2.261 & 2.670 & 3.361 \\
\texttt{degree\_preserving\_random} & 2.123 & 2.564 & 3.382 \\
\end{longtable}
}

The \texttt{identity\_no\_trade} variant --- which encodes no trade
edges --- consistently outperforms all trade-based graph constructions
across all horizons and both GNN families. This pattern is robust across
the 20 seed draws (Std MAE ≈ 0.03 pp). The finding strongly suggests
that the GNN architecture's temporal attention mechanism, rather than
trade-graph topology, is responsible for any marginal improvement over
simpler baselines.

\subsection{5.3 Probabilistic Forecast
Accuracy}\label{probabilistic-forecast-accuracy}

\textbf{Table 5: Probabilistic Results --- CRPS by Model and Horizon}
\emph{(lower is better)}

{\def\LTcaptype{none} % do not increment counter
\begin{longtable}[]{@{}
  >{\raggedright\arraybackslash}p{(\linewidth - 8\tabcolsep) * \real{0.1458}}
  >{\raggedright\arraybackslash}p{(\linewidth - 8\tabcolsep) * \real{0.2917}}
  >{\raggedright\arraybackslash}p{(\linewidth - 8\tabcolsep) * \real{0.1875}}
  >{\raggedright\arraybackslash}p{(\linewidth - 8\tabcolsep) * \real{0.1875}}
  >{\raggedright\arraybackslash}p{(\linewidth - 8\tabcolsep) * \real{0.1875}}@{}}
\toprule\noalign{}
\begin{minipage}[b]{\linewidth}\raggedright
Model
\end{minipage} & \begin{minipage}[b]{\linewidth}\raggedright
Graph Variant
\end{minipage} & \begin{minipage}[b]{\linewidth}\raggedright
H=1 CRPS
\end{minipage} & \begin{minipage}[b]{\linewidth}\raggedright
H=2 CRPS
\end{minipage} & \begin{minipage}[b]{\linewidth}\raggedright
H=4 CRPS
\end{minipage} \\
\midrule\noalign{}
\endhead
\bottomrule\noalign{}
\endlastfoot
\textbf{GCN} & \textbf{identity\_no\_trade} & \textbf{1.378} & 2.131 &
3.257 \\
MLP & --- & 1.444 & 2.300 & 3.553 \\
TCN & --- & 1.477 & 2.466 & 3.210 \\
Persistence & --- & 1.480 & 2.153 & 3.179 \\
ETS & --- & 1.481 & 2.153 & 3.179 \\
ARIMA & --- & 1.459 & 2.094 & 3.013 \\
\textbf{Temporal Graph} & \textbf{identity\_no\_trade} & 1.682 &
\textbf{2.010} & \textbf{2.451} \\
LSTM & --- & 1.842 & 2.129 & 2.581 \\
\end{longtable}
}

80\% prediction interval coverage for the best-performing models ranges
from 0.57--0.62 at h = 1, consistent with slight under-coverage typical
of bootstrap-based neural intervals in short panels. The Temporal Graph
with identity graph achieves 0.576 coverage at h = 2 (target: 0.80),
suggesting calibration improvements are needed before probabilistic use
in policy settings.

\subsection{5.4 Statistical
Significance}\label{statistical-significance}

\textbf{No graph model achieves statistically consistent superiority
over its best non-graph comparator across all 20 seeds.} Specifically:

\begin{itemize}
\tightlist
\item
  GCN with \texttt{directed\_trade} outperforms Ridge at h = 1 in 60\%
  of seeds (BH-adjusted \emph{p} \textless{} 0.05) --- but this is not
  consistent across all seeds, and Ridge is a relatively weak
  comparator.
\item
  Temporal Graph with \texttt{identity\_no\_trade} outperforms LSTM at h
  = 1 in 75\% of seeds, but only in 5\% of seeds at h = 2.
\item
  GCN with \texttt{identity\_no\_trade} outperforms Ridge in 80\% of
  seeds at h = 1, 65\% at h = 2, and only 5\% at h = 4.
\item
  No comparison clears the full 100\%-seed threshold against stronger
  comparators (ARIMA, TCN, Persistence).
\end{itemize}

These results indicate that the observed point-forecast gains of GNNs
are \textbf{not statistically robust} at conventional levels after
correcting for multiple comparisons and seed variability.

\begin{center}\rule{0.5\linewidth}{0.5pt}\end{center}

\section{6. Discussion}\label{discussion}

\subsection{6.1 The Identity Graph
Paradox}\label{the-identity-graph-paradox}

The most striking finding is that both GNN families perform best when
the trade graph is replaced by an identity matrix. This has two
interpretations:

\begin{enumerate}
\def\labelenumi{\arabic{enumi}.}
\tightlist
\item
  \textbf{Architecture effect:} The GCN and temporal attention
  mechanisms provide implicit regularisation or capacity advantages
  relative to simpler neural baselines, independent of cross-country
  propagation.
\item
  \textbf{Noisy edges:} Quarterly bilateral trade flows may be too
  coarse or noisy to carry reliable inflation-propagation signals at the
  frequencies studied. Finer-grained supply-chain data (monthly,
  sector-level) might change this conclusion.
\end{enumerate}

This pattern echoes findings in financial networks where random or null
graphs often match or exceed informative graph structures (Zhu et al.,
2021), suggesting that GNN over-parameterisation can mask the absence of
useful graph information.

\subsection{6.2 Horizon Dependence}\label{horizon-dependence}

The Temporal Graph's advantage grows with horizon: negligible at h = 1,
moderate at h = 2, and clearest at h = 4. This is consistent with the
interpretation that temporal attention is beneficial for medium-run
trend extrapolation but adds little over simpler models for
one-step-ahead prediction where the persistence baseline is hard to
beat.

\subsection{6.3 Implications for Policy}\label{implications-for-policy}

Our results suggest that central bank modellers should not assume that
trade-network GNNs improve upon well-calibrated univariate benchmarks
without rigorous, multi-seed ablation testing. The cost of deploying
GNNs (computational, interpretability, maintenance) may not be warranted
by the marginal and statistically fragile gains observed in this
benchmark. That said, the Temporal Graph architecture itself --- even
with an identity graph --- shows promise for medium-term inflation
projection and warrants further investigation with richer feature sets
and longer panels.

\subsection{6.4 Limitations}\label{limitations}

\begin{enumerate}
\def\labelenumi{\arabic{enumi}.}
\tightlist
\item
  \textbf{Panel scope:} 20 European economies; findings may not
  generalise to emerging markets or non-EU contexts with weaker
  trade-data quality.
\item
  \textbf{Revised vintages:} Models use revised, not real-time, data,
  which may overstate out-of-sample accuracy.
\item
  \textbf{Short initial window:} Training begins with only 14 quarters
  (2011Q2-- 2014Q4), limiting the stability of early model estimates.
\item
  \textbf{Pseudo-real-time only:} Information sets are lagged but not
  reconstructed from historical data releases; true real-time vintages
  would be required for central bank application.
\item
  \textbf{No causal claims:} All results pertain to predictive
  associations. Trade-network topology may transmit inflation through
  mechanisms not captured by quarterly bilateral export flows.
\end{enumerate}

\begin{center}\rule{0.5\linewidth}{0.5pt}\end{center}

\section{7. Conclusion}\label{conclusion}

We conduct a large-scale, fully reproducible prospective benchmark of
GNN inflation forecasting across 20 European economies at three
horizons, comparing 12 model families with 8 graph construction
strategies and 20 random seeds. Our principal findings are:

\begin{enumerate}
\def\labelenumi{\arabic{enumi}.}
\tightlist
\item
  \textbf{GNNs match or slightly exceed ARIMA and TCN at medium-to-long
  horizons} (h = 2, 4) in point forecasts, but the advantage is narrow
  and horizon-dependent.
\item
  \textbf{Trade-network topology does not explain GNN gains:} the
  identity (no-trade) graph consistently outperforms all trade-based
  graphs across both GNN families and all horizons.
\item
  \textbf{No performance difference is statistically robust} after
  multi-seed testing and FDR correction.
\end{enumerate}

These findings call for greater methodological rigour in the emerging
literature on GNN applications to macroeconomic forecasting: ablation
against null and identity graphs should be standard practice, and
multi-seed statistical testing is essential before claiming
network-based improvements.

Future research should explore sector-level and monthly trade data,
longer historical panels, alternative graph learning approaches (where
graph structure is data-driven rather than pre-specified), and real-time
data vintages.

\begin{center}\rule{0.5\linewidth}{0.5pt}\end{center}

\section{Data Availability}\label{data-availability}

All data are sourced from publicly available repositories:

\begin{itemize}
\tightlist
\item
  \textbf{CPI/HICP:} Eurostat HICP database; IMF International Financial
  Statistics
\item
  \textbf{Bilateral trade:} Eurostat Comext quarterly trade statistics
\item
  \textbf{Energy prices:} Eurostat energy component of HICP; World Bank
  commodity prices
\end{itemize}

Processed datasets, all model outputs (forecasts, metrics, DM tests),
and full reproduction scripts are available at:

\begin{quote}
\textbf{GitHub Repository:}
\url{https://github.com/mrayanasim09/MarketEquationDiscovery}

\textbf{SSRN Preprint (v1 baseline):}
\url{https://doi.org/10.2139/ssrn.7009041}
\end{quote}

All artefacts are SHA256-verified for exact reproducibility. The
benchmark can be re-run end-to-end using the provided
\texttt{reproduce.sh} script (estimated runtime: ≈ 12 hours on Apple
Silicon; longer on CPU-only hardware).

\section{Code Availability}\label{code-availability}

Full source code is publicly available at:
\url{https://github.com/mrayanasim09/MarketEquationDiscovery}

The repository includes:

\begin{itemize}
\tightlist
\item
  Benchmark protocol specification
  (\texttt{docs/research\_protocol\_v2\_1.md})
\item
  Benchmark execution engine
  (\texttt{src/models/run\_benchmark\_engine\_v2\_1.py})
\item
  Statistical analysis pipeline
  (\texttt{src/models/analyze\_v2\_1\_results.py})
\item
  Manuscript figure and table generation
  (\texttt{src/models/generate\_v2\_1\_manuscript.py})
\item
  SHA256-verified frozen results (\texttt{experiments/results/v2\_1/})
\end{itemize}

\section{Conflict of Interest}\label{conflict-of-interest}

The author declares no conflicts of interest.

\section{Funding}\label{funding}

This research received no external funding.

\section{Acknowledgements}\label{acknowledgements}

The author thanks the open-science community for the tools and datasets
that made this research possible. Bilateral trade data are sourced from
Eurostat Comext; macroeconomic indicators from the IMF International
Financial Statistics and World Bank World Development Indicators; CPI
data from Eurostat HICP.

\begin{center}\rule{0.5\linewidth}{0.5pt}\end{center}

\section{References}\label{references}

Bayoumi, T., Bui, T., \& Berkmen, P. (2023). \emph{Trade network
exposure and inflation dynamics.} IMF Working Paper WP/23/117.

Calvo, G. A., \& Reinhart, C. M. (2002). Fear of floating.
\emph{Quarterly Journal of Economics}, 117(2), 379--408.

Chen, Y., Li, M., \& Zhang, Z. (2023). Graph neural networks for
commodity price forecasting. \emph{Energy Economics}, 118, 106482.

Cheng, D., \& Zhu, J. (2022). Financial contagion detection via
spatio-temporal graph networks. \emph{Journal of Financial Stability},
63, 101073.

Diebold, F. X., \& Mariano, R. S. (1995). Comparing predictive accuracy.
\emph{Journal of Business \& Economic Statistics}, 13(3), 253--263.

Faust, J., \& Wright, J. H. (2013). Forecasting inflation. In G. Elliott
\& A. Timmermann (Eds.), \emph{Handbook of Economic Forecasting} (Vol.
2, pp.~2--56). Elsevier.

Forbes, K. J., \& Warnock, F. E. (2012). Capital flow waves: Surges,
stops, flight, and retrenchment. \emph{Journal of International
Economics}, 88(2), 235--251.

García-Martos, C., Rodríguez, J., \& Sánchez, M. J. (2015). Forecasting
electricity prices and their volatility using Unobserved Components.
\emph{Energy Economics}, 43, 218--228.

Harvey, D., Leybourne, S., \& Newbold, P. (1997). Testing the equality
of prediction mean squared errors. \emph{International Journal of
Forecasting}, 13(2), 281--291.

Hewamalage, H., Bergmeir, C., \& Bandara, K. (2021). Recurrent neural
networks for time series forecasting: Current status and future
directions. \emph{International Journal of Forecasting}, 37(1),
388--427.

Kipf, T. N., \& Welling, M. (2017). Semi-supervised classification with
graph convolutional networks. \emph{International Conference on Learning
Representations}.

Li, Y., Yu, R., Shahabi, C., \& Liu, Y. (2018). Diffusion convolutional
recurrent neural network: Data-driven traffic forecasting.
\emph{International Conference on Learning Representations}.

Makridakis, S., Spiliotis, E., \& Assimakopoulos, V. (2018). Statistical
and machine learning forecasting methods: Concerns and ways forward.
\emph{PLOS ONE}, 13(3), e0194889.

Medeiros, M. C., Vasconcelos, G. F. R., Veiga, Á., \& Zilberman, E.
(2021). Forecasting inflation in a data-rich environment: The benefits
of machine learning methods. \emph{Journal of Business \& Economic
Statistics}, 39(1), 98--119.

Miranda-Agrippino, S., \& Rey, H. (2021). The global financial cycle. In
G. Gopinath, E. Helpman, \& K. Rogoff (Eds.), \emph{Handbook of
International Economics} (Vol. 6, pp.~1--43). Elsevier.

Stock, J. H., \& Watson, M. W. (2007). Why has U.S. inflation become
harder to forecast? \emph{Journal of Money, Credit and Banking}, 39(s1),
3--33.

Wang, X., Chen, T., \& Li, H. (2024). Spatio-temporal graph networks for
regional GDP forecasting. \emph{Regional Science and Urban Economics},
105, 103990.

Wu, Z., Pan, S., Chen, F., Long, G., Zhang, C., \& Yu, P. S. (2019). A
comprehensive study on spatial-temporal graph neural networks.
\emph{IEEE Transactions on Neural Networks and Learning Systems}, 32(2),
527--540.

Zhu, L., Qiu, D., Ergu, D., Ying, C., \& Liu, K. (2021). A study on
predicting loan default based on the random forest algorithm.
\emph{Procedia Computer Science}, 176, 2337--2346.

\end{document}
